% ======================================================================
% Section 4 -- METHOD (~2 page budget). Direct register, no em dashes.
% ======================================================================
\section{Method: One Game, Two Outputs}
\label{sec:framework}

\method plays one cooperative game and reads two answers from it. The
players are the cells of the input complex. The value of a coalition is
the model's output when only those cells participate. The game returns
\emph{attributions} (per-cell Shapley values, \cref{eq:shapley}) and an
\emph{explanation} (the smallest sufficient sub-complex). The same
template, with neighborhoods as players and validation accuracy as the
value, explains \emph{performance} (\cref{sec:framework-game2}).

\subsection{The Participation Game}
\label{sec:framework-game1}

To explain a prediction $f(\mathcal{C})$ for class $y$, the players are
the cells $P \subseteq \mathcal{X}$ of the input, by default all cells,
with $n = |P|$. A coalition $S$ is scored on the \emph{induced
sub-complex} $\mathcal{C}[S]$: cells outside $S$ are detached, meaning
their features are removed, every neighborhood matrix is recomputed
from the restricted incidences, and readout pooling is limited to $S$:
\begin{equation}
\label{eq:participation-game}
v(S) \;=\;
  f_{y}\bigl(\mathcal{C}[S]\bigr) - f_{y}\bigl(\mathcal{C}[\emptyset]\bigr).
\end{equation}
Detachment is what absence means for a cell. Zeroing features instead
leaves the cell in place, still passing messages and still contributing
bias terms to pooling, so a zeroed game scores a different
counterfactual (quantified in \cref{app:experimental}). Detachment is
also cell-native: removing a face removes the rank-2 player itself.

\paragraph{Output 1: attributions.}
Shapley values $\phi_\sigma(v)$ on \cref{eq:participation-game} give
signed per-cell credit with the axioms of \cref{sec:prelim-shapley}.
Efficiency decomposes the prediction over cells, and the null-player
axiom detects models that ignore their higher-order structure
(\cref{sec:predictions}). Small complexes are enumerated exactly.
Larger ones use $T$ permutation passes, which cost $T \cdot n$ game
evaluations and are unbiased \citep{castro2009polynomial}.

\paragraph{Output 2: explanations.}
Motif recovery asks for a sub-complex, so \method extracts one
directly: the budget-$b$ coalition
$S^{*}_{b} \approx \argmax_{|S| = b} v(S)$, found by greedy pruning.
We start from $P$, repeatedly detach the cell whose removal least
reduces $v$, and record $S$ at every requested budget. One pass costs
at most $(n^{2}+n)/2$ evaluations of one forward pass each and yields
the coalition at all budgets, since pruning is nested. Pruning from
the full complex keeps evaluations near the training distribution.
Growing coalitions from single cells scores far outside it and fails
empirically (\cref{app:experimental}).

\paragraph{Relation to graph explainers.}
A graph is a complex of rank $\le 1$, where
\cref{eq:participation-game} reduces to induced-subgraph node removal.
Output 2 then solves the same search problem as SubgraphX
\citep{yuan2021subgraphx}, with two substitutions: detachment replaces
feature zero-filling, and exhaustive one-step pruning, at one forward
pass per scored coalition, replaces Monte-Carlo tree search whose
reward samples ${\sim}10^{2}$ context coalitions per node it scores.
Both substitutions are ablated in \cref{sec:predictions}. GNNExplainer
\citep{ying2019gnnexplainer} and PGExplainer
\citep{luo2020parameterized} optimize continuous relaxed masks and
return no game or axioms. All three are defined only at rank $\le 1$.

\subsection{Neighborhoods Explain Performance}
\label{sec:framework-game2}
\label{sec:method-pipeline}

The same template explains performance. The players are the model's
neighborhood-indexed components: a GCCN has one message-passing route
per neighborhood and HOPSE has one encoding block per neighborhood, so
switching a player off removes that neighborhood. We propose two value
functions. The \emph{retraining value} $v_{\mathrm{retrain}}(S)$ is
the validation accuracy of the model trained from scratch with the
neighborhoods in $S$; it is exact, and each evaluation costs a full
training run. The \emph{masked value} $v_{\mathrm{mask}}(S)$ is the
validation accuracy of a briefly trained full model evaluated with the
neighborhoods outside $S$ removed; each evaluation costs one forward
pass over the validation set.

We turn this game into an architecture selection procedure:
\begin{recipebox}
  \item We train the full model, with all candidate neighborhoods, for
    a small fixed number of epochs
    ($\sim$\unfrozen{\ResStemEpochs}).
  \item We repeatedly remove the neighborhood whose removal costs the
    least masked value, while that cost stays below the validation
    noise floor. This yields one candidate architecture per size.
  \item We train $B$ of these candidates to completion and keep the
    one with the best validation accuracy.
\end{recipebox}

The stop threshold is the binomial noise floor $\sqrt{p(1-p)/m}$ of a
validation set with $m$ examples, known before any training, so the
selected size is not a hyperparameter. The budget $B$ trades cost
against reliability, from $\sim$\unfrozen{\ResLadderCostBOne} training
runs at $B{=}1$ to $\sim$\unfrozen{\ResLadderCostBTwo} at $B{=}2$; a
full sweep trains every size. Two rules protect the procedure. The
masked value only orders candidates, and every accept or reject
decision is validation of fully trained models. Elimination walks
downward from the full model, the regime where the masked value is
most faithful. Details and supporting evidence are in
\cref{app:selection}.

\subsection{Exactness}
\label{sec:framework-exact}
\label{sec:framework-theory}

The neighborhood game has $n = \unfrozen{\ResNumNeighborhoods}$
players and $2^{n} = 2048$ coalitions, so \cref{eq:shapley,eq:interaction}
are evaluated exactly. Prediction games on small complexes likewise
enumerate; larger instances use sampled passes. Against brute force,
the implementation agrees to \unfrozen{\ResExactnessMaxErr}, which is
floating-point exactness. Exact axioms double as executable software
checks; the bugs they caught are documented in
\cref{app:experimental}.

\begin{theorybox}
\begin{propositionT}[Non-identifiability of substitutes; statement placeholder]
% (no \label: unnumbered T-statement; refer to it as "Proposition~T1")
% TODO(researcher): finalize statement + proof.
Let $N_a, N_b \in \mathcal{N}$ be neighborhoods that are functionally
substitutable for a trained model $f_\theta$, i.e., there exist
parameters realizing the same function with either route active. Then
the masked value computed on any single trained model cannot
distinguish the attribution profile of $(N_a, N_b)$ from that of a
model committed to either substitute. \emph{Exact statement TBD by
researcher.}
\end{propositionT}
\end{theorybox}

Proposition~T1 states which game answers which question: the masked
value identifies the trained model's realized division of labor, and
claims about what is realizable under retraining require the
retraining value. The selection procedure operationalizes this by
letting validation over trained models decide.
